\section{GCD, LCM y Euclides extendido}
\whenuse{Divisibilidad, simplificar fracciones, resolver $ax+by=gcd(a,b)$ o
construir inversos cuando los n\'umeros son coprimos.}
\complexity{$O(\log\min(|a|,|b|))$ tiempo; memoria logar\'itmica por el stack en
la versi\'on recursiva.}

\begin{lstlisting}[style=cpp]
optional<ll> checked_lcm(ll a, ll b) {
    if (a == 0 || b == 0) return 0;
    if (a == numeric_limits<ll>::min() || b == numeric_limits<ll>::min())
        return nullopt;
    ll divisor = gcd(a, b);
    i128 value = (i128)(a / divisor) * b;
    if (value < 0) value = -value;
    if (value > numeric_limits<ll>::max()) return nullopt;
    return (ll)value;
}

ll extended_gcd(ll a, ll b, ll &x, ll &y) {
    if (b == 0) { x = 1; y = 0; return a; }
    ll x1, y1;
    ll g = extended_gcd(b, a % b, x1, y1);
    x = y1;
    y = x1 - (a / b) * y1;
    return g;
}
\end{lstlisting}

\section{Aritm\'etica modular}
\whenuse{Potencias, divisiones e inversos bajo m\'odulo; el inverso por Fermat
requiere m\'odulo primo y denominador no divisible por \'el.}

\begin{lstlisting}[style=cpp]
ll mod_normalize(ll value, ll mod) {
    assert(mod > 0);
    value %= mod;
    if (value < 0) value += mod;
    return value;
}

ll mod_add(ll a, ll b, ll mod) {
    return ((i128)mod_normalize(a, mod) + mod_normalize(b, mod)) % mod;
}

ll mod_subtract(ll a, ll b, ll mod) {
    i128 result = (i128)mod_normalize(a, mod) - mod_normalize(b, mod);
    if (result < 0) result += mod;
    return (ll)result;
}

ll mod_multiply(ll a, ll b, ll mod) {
    return (i128)mod_normalize(a, mod) * mod_normalize(b, mod) % mod;
}

ll mod_power(ll base, unsigned long long exponent, ll mod) {
    ll result = 1 % mod;
    base = mod_normalize(base, mod);
    while (exponent) {
        if (exponent & 1ULL) result = mod_multiply(result, base, mod);
        base = mod_multiply(base, base, mod);
        exponent >>= 1;
    }
    return result;
}

optional<ll> mod_inverse(ll a, ll mod) {
    assert(mod > 0);
    ll x, y;
    ll g = extended_gcd(mod_normalize(a, mod), mod, x, y);
    if (g != 1) return nullopt;
    return mod_normalize(x, mod);
}

optional<ll> mod_divide(ll a, ll b, ll mod) {
    auto inverse = mod_inverse(b, mod);
    if (!inverse) return nullopt;
    return mod_multiply(a, *inverse, mod);
}
\end{lstlisting}

\complexity{$O(\log e)$ para potencia; $O(\log m)$ para inverso; memoria
$O(1)$ salvo el stack de Euclides.}

\section{Teorema chino del resto para dos congruencias}
Te permite saber \textbf{qu\'e n\'umeros cumplen las dos condiciones al mismo
tiempo} y \textbf{cada cu\'anto se repiten}.

\whenuse{Combinar condiciones $x\equiv a\pmod m$; esta versi\'on tambi\'en detecta
si dos congruencias con m\'odulos no coprimos son incompatibles.}
\complexity{$O(\log\min(m_1,m_2))$ tiempo por combinaci\'on; memoria logar\'itmica
por Euclides extendido.}

Admite m\'odulos no coprimos. Devuelve la soluci\'on normalizada y el m\'odulo
m\'inimo, o \texttt{nullopt} si no hay soluci\'on o el m\'odulo no entra en
\texttt{long long}.

\begin{lstlisting}[style=cpp]
optional<pair<ll,ll>> crt_pair(ll a1, ll m1, ll a2, ll m2) {
    assert(m1 > 0 && m2 > 0);
    a1 = mod_normalize(a1, m1);
    a2 = mod_normalize(a2, m2);
    ll x, y;
    ll g = extended_gcd(m1, m2, x, y);
    ll difference = a2 - a1;
    if (difference % g != 0) return nullopt;

    i128 modulus = (i128)(m1 / g) * m2;
    if (modulus > numeric_limits<ll>::max()) return nullopt;
    ll reduced = m2 / g;
    i128 t = (i128)(difference / g) * x % reduced;
    if (t < 0) t += reduced;
    i128 answer = ((i128)a1 + (i128)m1 * t) % modulus;
    if (answer < 0) answer += modulus;
    return pair<ll,ll>{(ll)answer, (ll)modulus};
}
\end{lstlisting}
